% =============================================================================
%  trpc-preamble.tex — preamble ร่วมของเอกสารชุด "สัญญา tRPC" ของโครงการ ULMs
%
%  ใช้โดย:  trpc-guide.tex    (ฉบับที่ 1 — หลักการ tRPC + ความเสี่ยงใน repo นี้)
%           trpc-meeting.tex  (ฉบับที่ 2 — วาระประชุมเรียงตามลำดับความสำคัญ)
%
%  ทำไมไม่ใช้ ulms-preamble.tex ที่มีอยู่แล้ว
%  ----------------------------------------
%  ulms-preamble.tex ตั้งใจให้เป็น "โทนเดียว (monochrome)" สำหรับเอกสารที่ต้อง
%  ส่งอาจารย์/พิมพ์ขาวดำ แต่เอกสารสอนสองฉบับนี้ใช้ "สีเชิงความหมาย" เป็นเครื่องมือ
%  สอนหลัก (สีเดียวกัน = แนวคิดเดียวกัน ทั้งในเนื้อความ ในไดอะแกรม และในตาราง)
%  จึงต้องมี preamble ของตัวเอง — ไม่ใช่ก๊อปมาแก้สีทับ
%
%  เอกสารที่เรียกใช้ต้องทำสองอย่างหลัง % =============================================================================
%  trpc-preamble.tex — preamble ร่วมของเอกสารชุด "สัญญา tRPC" ของโครงการ ULMs
%
%  ใช้โดย:  trpc-guide.tex    (ฉบับที่ 1 — หลักการ tRPC + ความเสี่ยงใน repo นี้)
%           trpc-meeting.tex  (ฉบับที่ 2 — วาระประชุมเรียงตามลำดับความสำคัญ)
%
%  ทำไมไม่ใช้ ulms-preamble.tex ที่มีอยู่แล้ว
%  ----------------------------------------
%  ulms-preamble.tex ตั้งใจให้เป็น "โทนเดียว (monochrome)" สำหรับเอกสารที่ต้อง
%  ส่งอาจารย์/พิมพ์ขาวดำ แต่เอกสารสอนสองฉบับนี้ใช้ "สีเชิงความหมาย" เป็นเครื่องมือ
%  สอนหลัก (สีเดียวกัน = แนวคิดเดียวกัน ทั้งในเนื้อความ ในไดอะแกรม และในตาราง)
%  จึงต้องมี preamble ของตัวเอง — ไม่ใช่ก๊อปมาแก้สีทับ
%
%  เอกสารที่เรียกใช้ต้องทำสองอย่างหลัง % =============================================================================
%  trpc-preamble.tex — preamble ร่วมของเอกสารชุด "สัญญา tRPC" ของโครงการ ULMs
%
%  ใช้โดย:  trpc-guide.tex    (ฉบับที่ 1 — หลักการ tRPC + ความเสี่ยงใน repo นี้)
%           trpc-meeting.tex  (ฉบับที่ 2 — วาระประชุมเรียงตามลำดับความสำคัญ)
%
%  ทำไมไม่ใช้ ulms-preamble.tex ที่มีอยู่แล้ว
%  ----------------------------------------
%  ulms-preamble.tex ตั้งใจให้เป็น "โทนเดียว (monochrome)" สำหรับเอกสารที่ต้อง
%  ส่งอาจารย์/พิมพ์ขาวดำ แต่เอกสารสอนสองฉบับนี้ใช้ "สีเชิงความหมาย" เป็นเครื่องมือ
%  สอนหลัก (สีเดียวกัน = แนวคิดเดียวกัน ทั้งในเนื้อความ ในไดอะแกรม และในตาราง)
%  จึงต้องมี preamble ของตัวเอง — ไม่ใช่ก๊อปมาแก้สีทับ
%
%  เอกสารที่เรียกใช้ต้องทำสองอย่างหลัง % =============================================================================
%  trpc-preamble.tex — preamble ร่วมของเอกสารชุด "สัญญา tRPC" ของโครงการ ULMs
%
%  ใช้โดย:  trpc-guide.tex    (ฉบับที่ 1 — หลักการ tRPC + ความเสี่ยงใน repo นี้)
%           trpc-meeting.tex  (ฉบับที่ 2 — วาระประชุมเรียงตามลำดับความสำคัญ)
%
%  ทำไมไม่ใช้ ulms-preamble.tex ที่มีอยู่แล้ว
%  ----------------------------------------
%  ulms-preamble.tex ตั้งใจให้เป็น "โทนเดียว (monochrome)" สำหรับเอกสารที่ต้อง
%  ส่งอาจารย์/พิมพ์ขาวดำ แต่เอกสารสอนสองฉบับนี้ใช้ "สีเชิงความหมาย" เป็นเครื่องมือ
%  สอนหลัก (สีเดียวกัน = แนวคิดเดียวกัน ทั้งในเนื้อความ ในไดอะแกรม และในตาราง)
%  จึงต้องมี preamble ของตัวเอง — ไม่ใช่ก๊อปมาแก้สีทับ
%
%  เอกสารที่เรียกใช้ต้องทำสองอย่างหลัง \input{trpc-preamble}:
%    1) \hypersetup{pdftitle={...}}
%    2) \begin{document}
%
%  Build:  cd docs && latexmk trpc-guide.tex     (.latexmkrc บังคับ xelatex ให้แล้ว)
%          PDF ออกที่ docs/build/ แล้วค่อยคัดลอกไป docs/report/
% =============================================================================

\documentclass[11pt,a4paper]{article}

% ---------------------------------------------------------------------
%  FONTS
%
%  ต้องใช้ XeLaTeX เท่านั้น — pdflatex เรียงพิมพ์ภาษาไทยไม่ได้เลย (ไม่ใช่ "ได้แต่ไม่สวย")
%
%  ใช้ตระกูล TLWG ที่มีทั้งอักษรไทยและละตินอยู่ในไฟล์เดียว จงใจเลี่ยงการจับคู่
%  ฟอนต์ละตินกับฟอนต์ไทยล้วนผ่าน ucharclasses เพราะ ucharclasses สลับฟอนต์ตาม
%  "บล็อกยูนิโคด" ของตัวอักษร โดยไม่รู้จักขอบเขตของ TeX group → ข้อความละตินที่
%  ตามหลัง \textbf{...} จะยังค้างอยู่ในฟอนต์ไทยซึ่งไม่มีสัญลักษณ์ละติน แล้ว
%  "หายไปจาก PDF เงียบ ๆ" โดยไม่มี warning
% ---------------------------------------------------------------------
\usepackage{fontspec}

\setmainfont{Laksaman}[Ligatures=TeX]          % ไทย+ละติน — เนื้อความ
\setsansfont{Garuda}[Ligatures=TeX]            % ไทย+ละติน หนักกว่า — หัวข้อ/ป้ายในไดอะแกรม
\setmonofont{DejaVu Sans Mono}[Scale=0.84]     % 0.84 ทำให้ x-height ใกล้เคียง TLWG

% --- การตัดบรรทัดภาษาไทย ---
% ภาษาไทยเขียนติดกันไม่มีช่องว่าง TeX จึงหาจุดตัดบรรทัดเองไม่ได้ และจะปล่อยให้
% บรรทัดล้นออกนอกหน้ากระดาษ สองบรรทัดนี้ยกงานให้ ICU ตัดคำแบบพจนานุกรม
% ค่า stretch ต้องมากกว่าศูนย์อย่างมีนัยสำคัญ ไม่งั้น TeX แทบไม่มีอิสระในการจัด
% บรรทัดไทยและจะออกมาเป็น overfull box แทนที่จะยืด/หดช่องไฟ
\XeTeXlinebreaklocale "th"
\XeTeXlinebreakskip = 0pt plus 0.15em minus 0.02em
\emergencystretch=2.2em

% หัวข้อ ป้ายในไดอะแกรม หัวตาราง header/footer ใช้ sans — เนื้อความไม่ใช้
% นี่คือสิ่งที่แยก "โครงสร้างสำหรับกวาดสายตา" ออกจาก "เนื้อหาสำหรับอ่าน"
\newcommand{\sansall}{\sffamily}

% ---------------------------------------------------------------------
%  LAYOUT
% ---------------------------------------------------------------------
\usepackage[a4paper,margin=2.3cm,top=2.5cm,bottom=2.5cm,headsep=12pt]{geometry}
\usepackage{amsmath,amssymb}
\usepackage{graphicx}
\graphicspath{{figures/}}
\usepackage{booktabs}
\usepackage{array}
\usepackage{longtable}
\usepackage{multirow}
\usepackage{xcolor}
\usepackage{tikz}
\usetikzlibrary{arrows.meta,positioning,calc,shapes.geometric,shapes.misc,
                fit,backgrounds,decorations.pathmorphing}
\usepackage[most]{tcolorbox}
\usepackage[font=small,labelfont=bf,width=0.95\textwidth]{caption}
\usepackage{fancyhdr}
\usepackage{enumitem}
\usepackage{titlesec}
\usepackage{float}        % figure[H] — ตรึงรูปไว้ข้างเนื้อความที่อธิบายมัน
\usepackage{microtype}
\usepackage{listings}
\usepackage[colorlinks=true,allcolors=ctr,breaklinks=true,
            bookmarksnumbered=true]{hyperref}

% ภาษาไทยต้องการระยะบรรทัดมากกว่าละติน สระบนซ้อนวรรณยุกต์ และมีสระล่างด้วย
% ที่ระยะปกติจะชนกันข้ามบรรทัด — 1.28 คือค่าที่ทดสอบแล้ว
\linespread{1.28}
\setlength{\parskip}{0.35em}
\setlength{\parindent}{0pt}

% ---------------------------------------------------------------------
%  สี — ระบบสีเชิงความหมาย (semantic colour)
%
%  ฐานสีคือชุด Okabe–Ito ซึ่งปลอดภัยสำหรับผู้ที่มองสีบางคู่ไม่แยก (colour-blind safe)
%
%  สีสี่สีหลักถูก "ตั้งชื่อกลาง ๆ" ไว้ที่นี่ แล้วให้แต่ละเอกสารประกาศความหมายของ
%  ตัวเองในหน้า "เอกสารนี้อ่านอย่างไร":
%
%    trpc-guide.tex   → 4 ชั้นของระบบ:   สัญญา / backend / frontend / เครือข่าย
%    trpc-meeting.tex → 4 ระดับเร่งด่วน: P0 / P1 / P2 / P3
%
%  สองเอกสารใช้คนละความหมายได้ เพราะแต่ละฉบับ "ประกาศระบบสีของตัวเองไว้หน้าแรก"
%  และคงเส้นคงวาภายในเล่ม — สิ่งที่ห้ามคือเปลี่ยนความหมายกลางเล่ม
% ---------------------------------------------------------------------
\definecolor{ctr}{HTML}{0072B2}      % ฟ้า   — สัญญา/schema/type   · หรือ P2
\definecolor{beo}{HTML}{D55E00}      % ส้ม   — ฝั่ง backend         · หรือ P1
\definecolor{feg}{HTML}{009E73}      % เขียว — ฝั่ง frontend        · หรือ P3
\definecolor{ops}{HTML}{7B5EA7}      % ม่วง  — เครือข่าย/รันไทม์/ops

\definecolor{ink}{HTML}{1A2733}      % เข้ม — หัวข้อรอง
\definecolor{muted}{HTML}{5B6B7A}    % ข้อความรอง เส้นคั่น ป้ายกำกับ
\definecolor{rule}{HTML}{C9D2DB}     % เส้นบาง ขอบไดอะแกรม
\definecolor{wash}{HTML}{F1F6FA}     % พื้นกล่องฟ้า/กลาง
\definecolor{warnwash}{HTML}{FDF3E7} % พื้นกล่องส้ม
\definecolor{goodwash}{HTML}{E9F4F0} % พื้นกล่องเขียว
\definecolor{badwash}{HTML}{FBEDED}  % พื้นกล่องแดง
\definecolor{bad}{HTML}{B3261E}      % แดง — ผิด/อันตราย · หรือ P0
\definecolor{codebg}{HTML}{F7F9FB}   % พื้นบล็อกโค้ด

% ---------------------------------------------------------------------
%  HEADINGS
%  \part ใช้เลขอารบิก เพื่อให้อ้างอิงข้ามเล่มอ่านว่า "ภาค 3" ไม่ใช่ "ภาค III"
% ---------------------------------------------------------------------
\renewcommand{\thepart}{\arabic{part}}
\titleformat{\part}[display]
  {\sansall\huge\bfseries\color{ctr}}
  {\normalsize\color{muted}ภาค \thepart}{0.4em}{}
\titlespacing*{\part}{0pt}{0pt}{1.5em}

\titleformat{\section}
  {\sansall\Large\bfseries\color{ctr!88!black}}{\thesection}{0.7em}{}
\titleformat{\subsection}
  {\sansall\large\bfseries\color{ink}}{\thesubsection}{0.6em}{}
\titleformat{\subsubsection}
  {\sansall\normalsize\bfseries\color{muted}}{\thesubsubsection}{0.55em}{}
\titlespacing*{\section}{0pt}{1.9em}{0.7em}
\titlespacing*{\subsection}{0pt}{1.3em}{0.45em}
\titlespacing*{\subsubsection}{0pt}{1.0em}{0.35em}

% หัวกระดาษ: ชื่อเอกสารคงที่ทางซ้าย ชื่อหัวข้อปัจจุบันทางขวา ผู้อ่านจะรู้ตลอดว่าอยู่ตรงไหน
% \docshort ถูกกำหนดในไฟล์เอกสารแต่ละฉบับ
\providecommand{\docshort}{}
\pagestyle{fancy}
\fancyhf{}
\renewcommand{\headrulewidth}{0.4pt}
\fancyhead[L]{\footnotesize\sansall\color{muted}\docshort}
\fancyhead[R]{\footnotesize\sansall\color{muted}\leftmark}
\fancyfoot[C]{\footnotesize\color{muted}\thepage}
\renewcommand{\sectionmark}[1]{\markboth{#1}{}}

% รายการแบบกระชับ — ค่าเริ่มต้นของ LaTeX ดูหลวมเมื่ออยู่กับ \parskip
\setlist[itemize]{itemsep=1pt,topsep=3pt,parsep=0pt,leftmargin=1.4em}
\setlist[enumerate]{itemsep=1pt,topsep=3pt,parsep=0pt,leftmargin=1.6em}
\setlist[description]{itemsep=2pt,topsep=3pt,parsep=0pt}

% ---------------------------------------------------------------------
%  กล่องข้อความ 4 แบบ ความหมายตายตัว ห้ามคิดแบบที่ห้ากลางเล่ม
%
%  colbacktitle ต้องตั้งให้เท่ากับ colback มิฉะนั้น tcolorbox จะใช้ colframe
%  (สีเข้ม) เป็นพื้นหลังของหัวกล่อง แล้วตัวอักษรหัวกล่องสีเข้มจะอ่านไม่ออก
% ---------------------------------------------------------------------
\tcbset{guidebox/.style={
  enhanced, breakable, boxrule=0pt, leftrule=2.5pt, arc=1pt,
  left=8pt, right=8pt, top=6pt, bottom=6pt,
  fonttitle=\sansall\bfseries\small, titlerule=0pt,
  toptitle=1pt, bottomtitle=1pt}}

% ฟ้า — ใจความที่ต้องจำ / ประโยคที่ร้อยหลายบทเข้าด้วยกัน
\newtcolorbox{keybox}[1][]{guidebox, colback=wash, colframe=ctr,
  colbacktitle=wash, coltitle=ctr!85!black, #1}
% ส้ม — กับดัก หรือข้อสมมติที่ถ้าลืมแล้วข้อสรุปจะผิด
\newtcolorbox{warnbox}[1][]{guidebox, colback=warnwash, colframe=beo,
  colbacktitle=warnwash, coltitle=beo!80!black, #1}
% เขียว — สิ่งที่ตรวจสอบกับของจริงในรีโปแล้ว
\newtcolorbox{goodbox}[1][]{guidebox, colback=goodwash, colframe=feg,
  colbacktitle=goodwash, coltitle=feg!65!black, #1}
% แดง — วิธีที่ผิดจริง ๆ อย่าทำ
\newtcolorbox{badbox}[1][]{guidebox, colback=badwash, colframe=bad,
  colbacktitle=badwash, coltitle=bad, #1}

% ---------------------------------------------------------------------
%  SEMANTIC MARKUP
% ---------------------------------------------------------------------
% ศัพท์อังกฤษที่แทรกในประโยคภาษาไทย ฟอนต์เนื้อความครอบคลุมละตินอยู่แล้วจึงไม่ต้อง
% สลับฟอนต์ — เก็บแมโครไว้เป็น "เครื่องหมายเชิงความหมาย" เผื่อวันหลังอยากเปลี่ยน
% การแสดงผลทั้งเล่มพร้อมกัน
\newcommand{\en}[1]{#1}
% ชื่อไฟล์ ฟังก์ชัน หรือค่าที่ปรากฏในโค้ดจริง
% ไม่บังคับขนาดตัวอักษร ปล่อยให้รับขนาดจากบริบท (สำคัญในตารางที่ใช้ \footnotesize
% มิฉะนั้นโค้ดจะกลายเป็นตัวใหญ่กว่าข้อความรอบข้างแล้วดันคอลัมน์ล้น)
% ฟอนต์ mono ถูกย่อไว้ 0.84 แล้วใน \setmonofont จึงพอดีกับเนื้อความอยู่แล้ว
\newcommand{\code}[1]{{\ttfamily #1}}

% ชื่อไฟล์ที่มีอักษรไทยอยู่ในชื่อ — ห้ามใช้ \code เพราะ DejaVu Sans Mono ไม่มี
% สัญลักษณ์ไทย ตัวอักษรจะหายไปจาก PDF โดยขึ้นแค่ warning "Missing character"
% ใช้ Garuda (sans ที่มีทั้งไทยและละติน) แทน
\newcommand{\fnth}[1]{{\sffamily\small #1}}

% ป้ายชื่อแนวคิด — สีเดียวกับที่ใช้ในไดอะแกรมและตาราง
\newcommand{\stg}[2]{\textcolor{#1}{\textbf{\en{#2}}}}

% ---------------------------------------------------------------------
%  ไวยากรณ์ของไดอะแกรม — นิยามครั้งเดียว ทุกรูปจะดูเป็นชุดเดียวกัน
%  หมายเหตุ: node ที่มีข้อความไทยต้องเผื่อ inner ysep มากกว่าไดอะแกรมภาษาละติน
%  เพราะสระบน/ล่างกินที่แนวตั้ง
% ---------------------------------------------------------------------
\tikzset{
  blk/.style={draw=rule, line width=0.6pt, rounded corners=2.5pt,
              fill=wash, minimum height=9mm, inner xsep=6pt, inner ysep=5pt,
              align=center, font=\sansall\footnotesize},
  stage/.style={blk, fill=ctr!8, draw=ctr!45},
  bestage/.style={blk, fill=beo!8, draw=beo!45},
  festage/.style={blk, fill=feg!8, draw=feg!45},
  opstage/.style={blk, fill=ops!8, draw=ops!45},
  ar/.style={-{Stealth[length=2.2mm,width=1.8mm]}, line width=0.7pt,
             draw=muted},
  lbl/.style={font=\sansall\scriptsize, text=muted, align=center},
  tag/.style={font=\sansall\scriptsize\bfseries, align=center},
}

% ---------------------------------------------------------------------
%  ตาราง
%  คอลัมน์แบบชิดซ้าย: ข้อความจัดชิดขอบสองข้าง (justified) ในคอลัมน์แคบจะเปิด
%  ช่องไฟกว้างมาก และแย่เป็นพิเศษกับภาษาไทย เพราะ ICU ตัดไทยเป็นก้อนที่ใหญ่กว่า
%  คำภาษาอังกฤษ → ใช้ L{} แทน p{} ในทุกคอลัมน์ที่มีข้อความยาว
% ---------------------------------------------------------------------
\newcolumntype{L}[1]{>{\raggedright\arraybackslash}p{#1}}
\captionsetup{skip=6pt}

% ---------------------------------------------------------------------
%  บล็อกโค้ด
%
%  ข้อจำกัดที่ต้องรู้: listings อ่านไฟล์แบบไบต์ ไม่เข้าใจ UTF-8 หลายไบต์
%  → "ห้ามใส่ภาษาไทยในบล็อกโค้ด" คอมเมนต์ในโค้ดทุกอันเป็นภาษาอังกฤษ
%    คำอธิบายภาษาไทยให้อยู่นอกบล็อกเสมอ
% ---------------------------------------------------------------------
\lstdefinelanguage{TypeScript}{
  keywords={await, async, break, case, catch, class, const, continue, default,
            delete, do, else, enum, export, extends, false, finally, for, from,
            function, if, implements, import, in, instanceof, interface, let,
            new, null, of, private, public, readonly, return, super, switch,
            this, throw, true, try, type, typeof, undefined, var, void, while,
            yield, satisfies, as},
  keywordstyle=\color{ctr}\bfseries,
  ndkeywords={string, number, boolean, Date, Promise, z, Record, Array},
  ndkeywordstyle=\color{ops},
  sensitive=true,
  comment=[l]{//},
  morecomment=[s]{/*}{*/},
  commentstyle=\color{muted}\itshape,
  stringstyle=\color{feg!70!black},
  morestring=[b]',
  morestring=[b]",
  morestring=[b]`,
}

\lstset{
  language=TypeScript,
  basicstyle=\ttfamily\scriptsize,
  backgroundcolor=\color{codebg},
  frame=leftline,
  framerule=2pt,
  rulecolor=\color{rule},
  framesep=7pt,
  xleftmargin=10pt,
  breaklines=true,
  breakatwhitespace=true,
  columns=fullflexible,
  keepspaces=true,
  showstringspaces=false,
  aboveskip=0.9em,
  belowskip=0.9em,
  literate={->}{{$\rightarrow$}}2 {=>}{{=>}}2,
}

% ชื่อ float ภาษาไทย — ถ้าไม่ตั้ง จะได้ "Figure 3" อยู่ใต้เนื้อความภาษาไทย
\renewcommand{\figurename}{รูปที่}
\renewcommand{\tablename}{ตารางที่}
\renewcommand{\refname}{เอกสารอ้างอิง}
\renewcommand{\contentsname}{สารบัญ}
% หมายเหตุ: ห้ามใส่ \color ใน \contentsname เพราะ \tableofcontents ส่งมันผ่าน
% \MakeUppercase ซึ่งจะเปลี่ยน ctr เป็น CTR แล้ว xcolor จะฟ้องว่าไม่รู้จักสีนั้น
% ถ้าจะจัดสไตล์หัวสารบัญ ให้ทำผ่าน \titleformat{\section} แทน
:
%    1) \hypersetup{pdftitle={...}}
%    2) \begin{document}
%
%  Build:  cd docs && latexmk trpc-guide.tex     (.latexmkrc บังคับ xelatex ให้แล้ว)
%          PDF ออกที่ docs/build/ แล้วค่อยคัดลอกไป docs/report/
% =============================================================================

\documentclass[11pt,a4paper]{article}

% ---------------------------------------------------------------------
%  FONTS
%
%  ต้องใช้ XeLaTeX เท่านั้น — pdflatex เรียงพิมพ์ภาษาไทยไม่ได้เลย (ไม่ใช่ "ได้แต่ไม่สวย")
%
%  ใช้ตระกูล TLWG ที่มีทั้งอักษรไทยและละตินอยู่ในไฟล์เดียว จงใจเลี่ยงการจับคู่
%  ฟอนต์ละตินกับฟอนต์ไทยล้วนผ่าน ucharclasses เพราะ ucharclasses สลับฟอนต์ตาม
%  "บล็อกยูนิโคด" ของตัวอักษร โดยไม่รู้จักขอบเขตของ TeX group → ข้อความละตินที่
%  ตามหลัง \textbf{...} จะยังค้างอยู่ในฟอนต์ไทยซึ่งไม่มีสัญลักษณ์ละติน แล้ว
%  "หายไปจาก PDF เงียบ ๆ" โดยไม่มี warning
% ---------------------------------------------------------------------
\usepackage{fontspec}

\setmainfont{Laksaman}[Ligatures=TeX]          % ไทย+ละติน — เนื้อความ
\setsansfont{Garuda}[Ligatures=TeX]            % ไทย+ละติน หนักกว่า — หัวข้อ/ป้ายในไดอะแกรม
\setmonofont{DejaVu Sans Mono}[Scale=0.84]     % 0.84 ทำให้ x-height ใกล้เคียง TLWG

% --- การตัดบรรทัดภาษาไทย ---
% ภาษาไทยเขียนติดกันไม่มีช่องว่าง TeX จึงหาจุดตัดบรรทัดเองไม่ได้ และจะปล่อยให้
% บรรทัดล้นออกนอกหน้ากระดาษ สองบรรทัดนี้ยกงานให้ ICU ตัดคำแบบพจนานุกรม
% ค่า stretch ต้องมากกว่าศูนย์อย่างมีนัยสำคัญ ไม่งั้น TeX แทบไม่มีอิสระในการจัด
% บรรทัดไทยและจะออกมาเป็น overfull box แทนที่จะยืด/หดช่องไฟ
\XeTeXlinebreaklocale "th"
\XeTeXlinebreakskip = 0pt plus 0.15em minus 0.02em
\emergencystretch=2.2em

% หัวข้อ ป้ายในไดอะแกรม หัวตาราง header/footer ใช้ sans — เนื้อความไม่ใช้
% นี่คือสิ่งที่แยก "โครงสร้างสำหรับกวาดสายตา" ออกจาก "เนื้อหาสำหรับอ่าน"
\newcommand{\sansall}{\sffamily}

% ---------------------------------------------------------------------
%  LAYOUT
% ---------------------------------------------------------------------
\usepackage[a4paper,margin=2.3cm,top=2.5cm,bottom=2.5cm,headsep=12pt]{geometry}
\usepackage{amsmath,amssymb}
\usepackage{graphicx}
\graphicspath{{figures/}}
\usepackage{booktabs}
\usepackage{array}
\usepackage{longtable}
\usepackage{multirow}
\usepackage{xcolor}
\usepackage{tikz}
\usetikzlibrary{arrows.meta,positioning,calc,shapes.geometric,shapes.misc,
                fit,backgrounds,decorations.pathmorphing}
\usepackage[most]{tcolorbox}
\usepackage[font=small,labelfont=bf,width=0.95\textwidth]{caption}
\usepackage{fancyhdr}
\usepackage{enumitem}
\usepackage{titlesec}
\usepackage{float}        % figure[H] — ตรึงรูปไว้ข้างเนื้อความที่อธิบายมัน
\usepackage{microtype}
\usepackage{listings}
\usepackage[colorlinks=true,allcolors=ctr,breaklinks=true,
            bookmarksnumbered=true]{hyperref}

% ภาษาไทยต้องการระยะบรรทัดมากกว่าละติน สระบนซ้อนวรรณยุกต์ และมีสระล่างด้วย
% ที่ระยะปกติจะชนกันข้ามบรรทัด — 1.28 คือค่าที่ทดสอบแล้ว
\linespread{1.28}
\setlength{\parskip}{0.35em}
\setlength{\parindent}{0pt}

% ---------------------------------------------------------------------
%  สี — ระบบสีเชิงความหมาย (semantic colour)
%
%  ฐานสีคือชุด Okabe–Ito ซึ่งปลอดภัยสำหรับผู้ที่มองสีบางคู่ไม่แยก (colour-blind safe)
%
%  สีสี่สีหลักถูก "ตั้งชื่อกลาง ๆ" ไว้ที่นี่ แล้วให้แต่ละเอกสารประกาศความหมายของ
%  ตัวเองในหน้า "เอกสารนี้อ่านอย่างไร":
%
%    trpc-guide.tex   → 4 ชั้นของระบบ:   สัญญา / backend / frontend / เครือข่าย
%    trpc-meeting.tex → 4 ระดับเร่งด่วน: P0 / P1 / P2 / P3
%
%  สองเอกสารใช้คนละความหมายได้ เพราะแต่ละฉบับ "ประกาศระบบสีของตัวเองไว้หน้าแรก"
%  และคงเส้นคงวาภายในเล่ม — สิ่งที่ห้ามคือเปลี่ยนความหมายกลางเล่ม
% ---------------------------------------------------------------------
\definecolor{ctr}{HTML}{0072B2}      % ฟ้า   — สัญญา/schema/type   · หรือ P2
\definecolor{beo}{HTML}{D55E00}      % ส้ม   — ฝั่ง backend         · หรือ P1
\definecolor{feg}{HTML}{009E73}      % เขียว — ฝั่ง frontend        · หรือ P3
\definecolor{ops}{HTML}{7B5EA7}      % ม่วง  — เครือข่าย/รันไทม์/ops

\definecolor{ink}{HTML}{1A2733}      % เข้ม — หัวข้อรอง
\definecolor{muted}{HTML}{5B6B7A}    % ข้อความรอง เส้นคั่น ป้ายกำกับ
\definecolor{rule}{HTML}{C9D2DB}     % เส้นบาง ขอบไดอะแกรม
\definecolor{wash}{HTML}{F1F6FA}     % พื้นกล่องฟ้า/กลาง
\definecolor{warnwash}{HTML}{FDF3E7} % พื้นกล่องส้ม
\definecolor{goodwash}{HTML}{E9F4F0} % พื้นกล่องเขียว
\definecolor{badwash}{HTML}{FBEDED}  % พื้นกล่องแดง
\definecolor{bad}{HTML}{B3261E}      % แดง — ผิด/อันตราย · หรือ P0
\definecolor{codebg}{HTML}{F7F9FB}   % พื้นบล็อกโค้ด

% ---------------------------------------------------------------------
%  HEADINGS
%  \part ใช้เลขอารบิก เพื่อให้อ้างอิงข้ามเล่มอ่านว่า "ภาค 3" ไม่ใช่ "ภาค III"
% ---------------------------------------------------------------------
\renewcommand{\thepart}{\arabic{part}}
\titleformat{\part}[display]
  {\sansall\huge\bfseries\color{ctr}}
  {\normalsize\color{muted}ภาค \thepart}{0.4em}{}
\titlespacing*{\part}{0pt}{0pt}{1.5em}

\titleformat{\section}
  {\sansall\Large\bfseries\color{ctr!88!black}}{\thesection}{0.7em}{}
\titleformat{\subsection}
  {\sansall\large\bfseries\color{ink}}{\thesubsection}{0.6em}{}
\titleformat{\subsubsection}
  {\sansall\normalsize\bfseries\color{muted}}{\thesubsubsection}{0.55em}{}
\titlespacing*{\section}{0pt}{1.9em}{0.7em}
\titlespacing*{\subsection}{0pt}{1.3em}{0.45em}
\titlespacing*{\subsubsection}{0pt}{1.0em}{0.35em}

% หัวกระดาษ: ชื่อเอกสารคงที่ทางซ้าย ชื่อหัวข้อปัจจุบันทางขวา ผู้อ่านจะรู้ตลอดว่าอยู่ตรงไหน
% \docshort ถูกกำหนดในไฟล์เอกสารแต่ละฉบับ
\providecommand{\docshort}{}
\pagestyle{fancy}
\fancyhf{}
\renewcommand{\headrulewidth}{0.4pt}
\fancyhead[L]{\footnotesize\sansall\color{muted}\docshort}
\fancyhead[R]{\footnotesize\sansall\color{muted}\leftmark}
\fancyfoot[C]{\footnotesize\color{muted}\thepage}
\renewcommand{\sectionmark}[1]{\markboth{#1}{}}

% รายการแบบกระชับ — ค่าเริ่มต้นของ LaTeX ดูหลวมเมื่ออยู่กับ \parskip
\setlist[itemize]{itemsep=1pt,topsep=3pt,parsep=0pt,leftmargin=1.4em}
\setlist[enumerate]{itemsep=1pt,topsep=3pt,parsep=0pt,leftmargin=1.6em}
\setlist[description]{itemsep=2pt,topsep=3pt,parsep=0pt}

% ---------------------------------------------------------------------
%  กล่องข้อความ 4 แบบ ความหมายตายตัว ห้ามคิดแบบที่ห้ากลางเล่ม
%
%  colbacktitle ต้องตั้งให้เท่ากับ colback มิฉะนั้น tcolorbox จะใช้ colframe
%  (สีเข้ม) เป็นพื้นหลังของหัวกล่อง แล้วตัวอักษรหัวกล่องสีเข้มจะอ่านไม่ออก
% ---------------------------------------------------------------------
\tcbset{guidebox/.style={
  enhanced, breakable, boxrule=0pt, leftrule=2.5pt, arc=1pt,
  left=8pt, right=8pt, top=6pt, bottom=6pt,
  fonttitle=\sansall\bfseries\small, titlerule=0pt,
  toptitle=1pt, bottomtitle=1pt}}

% ฟ้า — ใจความที่ต้องจำ / ประโยคที่ร้อยหลายบทเข้าด้วยกัน
\newtcolorbox{keybox}[1][]{guidebox, colback=wash, colframe=ctr,
  colbacktitle=wash, coltitle=ctr!85!black, #1}
% ส้ม — กับดัก หรือข้อสมมติที่ถ้าลืมแล้วข้อสรุปจะผิด
\newtcolorbox{warnbox}[1][]{guidebox, colback=warnwash, colframe=beo,
  colbacktitle=warnwash, coltitle=beo!80!black, #1}
% เขียว — สิ่งที่ตรวจสอบกับของจริงในรีโปแล้ว
\newtcolorbox{goodbox}[1][]{guidebox, colback=goodwash, colframe=feg,
  colbacktitle=goodwash, coltitle=feg!65!black, #1}
% แดง — วิธีที่ผิดจริง ๆ อย่าทำ
\newtcolorbox{badbox}[1][]{guidebox, colback=badwash, colframe=bad,
  colbacktitle=badwash, coltitle=bad, #1}

% ---------------------------------------------------------------------
%  SEMANTIC MARKUP
% ---------------------------------------------------------------------
% ศัพท์อังกฤษที่แทรกในประโยคภาษาไทย ฟอนต์เนื้อความครอบคลุมละตินอยู่แล้วจึงไม่ต้อง
% สลับฟอนต์ — เก็บแมโครไว้เป็น "เครื่องหมายเชิงความหมาย" เผื่อวันหลังอยากเปลี่ยน
% การแสดงผลทั้งเล่มพร้อมกัน
\newcommand{\en}[1]{#1}
% ชื่อไฟล์ ฟังก์ชัน หรือค่าที่ปรากฏในโค้ดจริง
% ไม่บังคับขนาดตัวอักษร ปล่อยให้รับขนาดจากบริบท (สำคัญในตารางที่ใช้ \footnotesize
% มิฉะนั้นโค้ดจะกลายเป็นตัวใหญ่กว่าข้อความรอบข้างแล้วดันคอลัมน์ล้น)
% ฟอนต์ mono ถูกย่อไว้ 0.84 แล้วใน \setmonofont จึงพอดีกับเนื้อความอยู่แล้ว
\newcommand{\code}[1]{{\ttfamily #1}}

% ชื่อไฟล์ที่มีอักษรไทยอยู่ในชื่อ — ห้ามใช้ \code เพราะ DejaVu Sans Mono ไม่มี
% สัญลักษณ์ไทย ตัวอักษรจะหายไปจาก PDF โดยขึ้นแค่ warning "Missing character"
% ใช้ Garuda (sans ที่มีทั้งไทยและละติน) แทน
\newcommand{\fnth}[1]{{\sffamily\small #1}}

% ป้ายชื่อแนวคิด — สีเดียวกับที่ใช้ในไดอะแกรมและตาราง
\newcommand{\stg}[2]{\textcolor{#1}{\textbf{\en{#2}}}}

% ---------------------------------------------------------------------
%  ไวยากรณ์ของไดอะแกรม — นิยามครั้งเดียว ทุกรูปจะดูเป็นชุดเดียวกัน
%  หมายเหตุ: node ที่มีข้อความไทยต้องเผื่อ inner ysep มากกว่าไดอะแกรมภาษาละติน
%  เพราะสระบน/ล่างกินที่แนวตั้ง
% ---------------------------------------------------------------------
\tikzset{
  blk/.style={draw=rule, line width=0.6pt, rounded corners=2.5pt,
              fill=wash, minimum height=9mm, inner xsep=6pt, inner ysep=5pt,
              align=center, font=\sansall\footnotesize},
  stage/.style={blk, fill=ctr!8, draw=ctr!45},
  bestage/.style={blk, fill=beo!8, draw=beo!45},
  festage/.style={blk, fill=feg!8, draw=feg!45},
  opstage/.style={blk, fill=ops!8, draw=ops!45},
  ar/.style={-{Stealth[length=2.2mm,width=1.8mm]}, line width=0.7pt,
             draw=muted},
  lbl/.style={font=\sansall\scriptsize, text=muted, align=center},
  tag/.style={font=\sansall\scriptsize\bfseries, align=center},
}

% ---------------------------------------------------------------------
%  ตาราง
%  คอลัมน์แบบชิดซ้าย: ข้อความจัดชิดขอบสองข้าง (justified) ในคอลัมน์แคบจะเปิด
%  ช่องไฟกว้างมาก และแย่เป็นพิเศษกับภาษาไทย เพราะ ICU ตัดไทยเป็นก้อนที่ใหญ่กว่า
%  คำภาษาอังกฤษ → ใช้ L{} แทน p{} ในทุกคอลัมน์ที่มีข้อความยาว
% ---------------------------------------------------------------------
\newcolumntype{L}[1]{>{\raggedright\arraybackslash}p{#1}}
\captionsetup{skip=6pt}

% ---------------------------------------------------------------------
%  บล็อกโค้ด
%
%  ข้อจำกัดที่ต้องรู้: listings อ่านไฟล์แบบไบต์ ไม่เข้าใจ UTF-8 หลายไบต์
%  → "ห้ามใส่ภาษาไทยในบล็อกโค้ด" คอมเมนต์ในโค้ดทุกอันเป็นภาษาอังกฤษ
%    คำอธิบายภาษาไทยให้อยู่นอกบล็อกเสมอ
% ---------------------------------------------------------------------
\lstdefinelanguage{TypeScript}{
  keywords={await, async, break, case, catch, class, const, continue, default,
            delete, do, else, enum, export, extends, false, finally, for, from,
            function, if, implements, import, in, instanceof, interface, let,
            new, null, of, private, public, readonly, return, super, switch,
            this, throw, true, try, type, typeof, undefined, var, void, while,
            yield, satisfies, as},
  keywordstyle=\color{ctr}\bfseries,
  ndkeywords={string, number, boolean, Date, Promise, z, Record, Array},
  ndkeywordstyle=\color{ops},
  sensitive=true,
  comment=[l]{//},
  morecomment=[s]{/*}{*/},
  commentstyle=\color{muted}\itshape,
  stringstyle=\color{feg!70!black},
  morestring=[b]',
  morestring=[b]",
  morestring=[b]`,
}

\lstset{
  language=TypeScript,
  basicstyle=\ttfamily\scriptsize,
  backgroundcolor=\color{codebg},
  frame=leftline,
  framerule=2pt,
  rulecolor=\color{rule},
  framesep=7pt,
  xleftmargin=10pt,
  breaklines=true,
  breakatwhitespace=true,
  columns=fullflexible,
  keepspaces=true,
  showstringspaces=false,
  aboveskip=0.9em,
  belowskip=0.9em,
  literate={->}{{$\rightarrow$}}2 {=>}{{=>}}2,
}

% ชื่อ float ภาษาไทย — ถ้าไม่ตั้ง จะได้ "Figure 3" อยู่ใต้เนื้อความภาษาไทย
\renewcommand{\figurename}{รูปที่}
\renewcommand{\tablename}{ตารางที่}
\renewcommand{\refname}{เอกสารอ้างอิง}
\renewcommand{\contentsname}{สารบัญ}
% หมายเหตุ: ห้ามใส่ \color ใน \contentsname เพราะ \tableofcontents ส่งมันผ่าน
% \MakeUppercase ซึ่งจะเปลี่ยน ctr เป็น CTR แล้ว xcolor จะฟ้องว่าไม่รู้จักสีนั้น
% ถ้าจะจัดสไตล์หัวสารบัญ ให้ทำผ่าน \titleformat{\section} แทน
:
%    1) \hypersetup{pdftitle={...}}
%    2) \begin{document}
%
%  Build:  cd docs && latexmk trpc-guide.tex     (.latexmkrc บังคับ xelatex ให้แล้ว)
%          PDF ออกที่ docs/build/ แล้วค่อยคัดลอกไป docs/report/
% =============================================================================

\documentclass[11pt,a4paper]{article}

% ---------------------------------------------------------------------
%  FONTS
%
%  ต้องใช้ XeLaTeX เท่านั้น — pdflatex เรียงพิมพ์ภาษาไทยไม่ได้เลย (ไม่ใช่ "ได้แต่ไม่สวย")
%
%  ใช้ตระกูล TLWG ที่มีทั้งอักษรไทยและละตินอยู่ในไฟล์เดียว จงใจเลี่ยงการจับคู่
%  ฟอนต์ละตินกับฟอนต์ไทยล้วนผ่าน ucharclasses เพราะ ucharclasses สลับฟอนต์ตาม
%  "บล็อกยูนิโคด" ของตัวอักษร โดยไม่รู้จักขอบเขตของ TeX group → ข้อความละตินที่
%  ตามหลัง \textbf{...} จะยังค้างอยู่ในฟอนต์ไทยซึ่งไม่มีสัญลักษณ์ละติน แล้ว
%  "หายไปจาก PDF เงียบ ๆ" โดยไม่มี warning
% ---------------------------------------------------------------------
\usepackage{fontspec}

\setmainfont{Laksaman}[Ligatures=TeX]          % ไทย+ละติน — เนื้อความ
\setsansfont{Garuda}[Ligatures=TeX]            % ไทย+ละติน หนักกว่า — หัวข้อ/ป้ายในไดอะแกรม
\setmonofont{DejaVu Sans Mono}[Scale=0.84]     % 0.84 ทำให้ x-height ใกล้เคียง TLWG

% --- การตัดบรรทัดภาษาไทย ---
% ภาษาไทยเขียนติดกันไม่มีช่องว่าง TeX จึงหาจุดตัดบรรทัดเองไม่ได้ และจะปล่อยให้
% บรรทัดล้นออกนอกหน้ากระดาษ สองบรรทัดนี้ยกงานให้ ICU ตัดคำแบบพจนานุกรม
% ค่า stretch ต้องมากกว่าศูนย์อย่างมีนัยสำคัญ ไม่งั้น TeX แทบไม่มีอิสระในการจัด
% บรรทัดไทยและจะออกมาเป็น overfull box แทนที่จะยืด/หดช่องไฟ
\XeTeXlinebreaklocale "th"
\XeTeXlinebreakskip = 0pt plus 0.15em minus 0.02em
\emergencystretch=2.2em

% หัวข้อ ป้ายในไดอะแกรม หัวตาราง header/footer ใช้ sans — เนื้อความไม่ใช้
% นี่คือสิ่งที่แยก "โครงสร้างสำหรับกวาดสายตา" ออกจาก "เนื้อหาสำหรับอ่าน"
\newcommand{\sansall}{\sffamily}

% ---------------------------------------------------------------------
%  LAYOUT
% ---------------------------------------------------------------------
\usepackage[a4paper,margin=2.3cm,top=2.5cm,bottom=2.5cm,headsep=12pt]{geometry}
\usepackage{amsmath,amssymb}
\usepackage{graphicx}
\graphicspath{{figures/}}
\usepackage{booktabs}
\usepackage{array}
\usepackage{longtable}
\usepackage{multirow}
\usepackage{xcolor}
\usepackage{tikz}
\usetikzlibrary{arrows.meta,positioning,calc,shapes.geometric,shapes.misc,
                fit,backgrounds,decorations.pathmorphing}
\usepackage[most]{tcolorbox}
\usepackage[font=small,labelfont=bf,width=0.95\textwidth]{caption}
\usepackage{fancyhdr}
\usepackage{enumitem}
\usepackage{titlesec}
\usepackage{float}        % figure[H] — ตรึงรูปไว้ข้างเนื้อความที่อธิบายมัน
\usepackage{microtype}
\usepackage{listings}
\usepackage[colorlinks=true,allcolors=ctr,breaklinks=true,
            bookmarksnumbered=true]{hyperref}

% ภาษาไทยต้องการระยะบรรทัดมากกว่าละติน สระบนซ้อนวรรณยุกต์ และมีสระล่างด้วย
% ที่ระยะปกติจะชนกันข้ามบรรทัด — 1.28 คือค่าที่ทดสอบแล้ว
\linespread{1.28}
\setlength{\parskip}{0.35em}
\setlength{\parindent}{0pt}

% ---------------------------------------------------------------------
%  สี — ระบบสีเชิงความหมาย (semantic colour)
%
%  ฐานสีคือชุด Okabe–Ito ซึ่งปลอดภัยสำหรับผู้ที่มองสีบางคู่ไม่แยก (colour-blind safe)
%
%  สีสี่สีหลักถูก "ตั้งชื่อกลาง ๆ" ไว้ที่นี่ แล้วให้แต่ละเอกสารประกาศความหมายของ
%  ตัวเองในหน้า "เอกสารนี้อ่านอย่างไร":
%
%    trpc-guide.tex   → 4 ชั้นของระบบ:   สัญญา / backend / frontend / เครือข่าย
%    trpc-meeting.tex → 4 ระดับเร่งด่วน: P0 / P1 / P2 / P3
%
%  สองเอกสารใช้คนละความหมายได้ เพราะแต่ละฉบับ "ประกาศระบบสีของตัวเองไว้หน้าแรก"
%  และคงเส้นคงวาภายในเล่ม — สิ่งที่ห้ามคือเปลี่ยนความหมายกลางเล่ม
% ---------------------------------------------------------------------
\definecolor{ctr}{HTML}{0072B2}      % ฟ้า   — สัญญา/schema/type   · หรือ P2
\definecolor{beo}{HTML}{D55E00}      % ส้ม   — ฝั่ง backend         · หรือ P1
\definecolor{feg}{HTML}{009E73}      % เขียว — ฝั่ง frontend        · หรือ P3
\definecolor{ops}{HTML}{7B5EA7}      % ม่วง  — เครือข่าย/รันไทม์/ops

\definecolor{ink}{HTML}{1A2733}      % เข้ม — หัวข้อรอง
\definecolor{muted}{HTML}{5B6B7A}    % ข้อความรอง เส้นคั่น ป้ายกำกับ
\definecolor{rule}{HTML}{C9D2DB}     % เส้นบาง ขอบไดอะแกรม
\definecolor{wash}{HTML}{F1F6FA}     % พื้นกล่องฟ้า/กลาง
\definecolor{warnwash}{HTML}{FDF3E7} % พื้นกล่องส้ม
\definecolor{goodwash}{HTML}{E9F4F0} % พื้นกล่องเขียว
\definecolor{badwash}{HTML}{FBEDED}  % พื้นกล่องแดง
\definecolor{bad}{HTML}{B3261E}      % แดง — ผิด/อันตราย · หรือ P0
\definecolor{codebg}{HTML}{F7F9FB}   % พื้นบล็อกโค้ด

% ---------------------------------------------------------------------
%  HEADINGS
%  \part ใช้เลขอารบิก เพื่อให้อ้างอิงข้ามเล่มอ่านว่า "ภาค 3" ไม่ใช่ "ภาค III"
% ---------------------------------------------------------------------
\renewcommand{\thepart}{\arabic{part}}
\titleformat{\part}[display]
  {\sansall\huge\bfseries\color{ctr}}
  {\normalsize\color{muted}ภาค \thepart}{0.4em}{}
\titlespacing*{\part}{0pt}{0pt}{1.5em}

\titleformat{\section}
  {\sansall\Large\bfseries\color{ctr!88!black}}{\thesection}{0.7em}{}
\titleformat{\subsection}
  {\sansall\large\bfseries\color{ink}}{\thesubsection}{0.6em}{}
\titleformat{\subsubsection}
  {\sansall\normalsize\bfseries\color{muted}}{\thesubsubsection}{0.55em}{}
\titlespacing*{\section}{0pt}{1.9em}{0.7em}
\titlespacing*{\subsection}{0pt}{1.3em}{0.45em}
\titlespacing*{\subsubsection}{0pt}{1.0em}{0.35em}

% หัวกระดาษ: ชื่อเอกสารคงที่ทางซ้าย ชื่อหัวข้อปัจจุบันทางขวา ผู้อ่านจะรู้ตลอดว่าอยู่ตรงไหน
% \docshort ถูกกำหนดในไฟล์เอกสารแต่ละฉบับ
\providecommand{\docshort}{}
\pagestyle{fancy}
\fancyhf{}
\renewcommand{\headrulewidth}{0.4pt}
\fancyhead[L]{\footnotesize\sansall\color{muted}\docshort}
\fancyhead[R]{\footnotesize\sansall\color{muted}\leftmark}
\fancyfoot[C]{\footnotesize\color{muted}\thepage}
\renewcommand{\sectionmark}[1]{\markboth{#1}{}}

% รายการแบบกระชับ — ค่าเริ่มต้นของ LaTeX ดูหลวมเมื่ออยู่กับ \parskip
\setlist[itemize]{itemsep=1pt,topsep=3pt,parsep=0pt,leftmargin=1.4em}
\setlist[enumerate]{itemsep=1pt,topsep=3pt,parsep=0pt,leftmargin=1.6em}
\setlist[description]{itemsep=2pt,topsep=3pt,parsep=0pt}

% ---------------------------------------------------------------------
%  กล่องข้อความ 4 แบบ ความหมายตายตัว ห้ามคิดแบบที่ห้ากลางเล่ม
%
%  colbacktitle ต้องตั้งให้เท่ากับ colback มิฉะนั้น tcolorbox จะใช้ colframe
%  (สีเข้ม) เป็นพื้นหลังของหัวกล่อง แล้วตัวอักษรหัวกล่องสีเข้มจะอ่านไม่ออก
% ---------------------------------------------------------------------
\tcbset{guidebox/.style={
  enhanced, breakable, boxrule=0pt, leftrule=2.5pt, arc=1pt,
  left=8pt, right=8pt, top=6pt, bottom=6pt,
  fonttitle=\sansall\bfseries\small, titlerule=0pt,
  toptitle=1pt, bottomtitle=1pt}}

% ฟ้า — ใจความที่ต้องจำ / ประโยคที่ร้อยหลายบทเข้าด้วยกัน
\newtcolorbox{keybox}[1][]{guidebox, colback=wash, colframe=ctr,
  colbacktitle=wash, coltitle=ctr!85!black, #1}
% ส้ม — กับดัก หรือข้อสมมติที่ถ้าลืมแล้วข้อสรุปจะผิด
\newtcolorbox{warnbox}[1][]{guidebox, colback=warnwash, colframe=beo,
  colbacktitle=warnwash, coltitle=beo!80!black, #1}
% เขียว — สิ่งที่ตรวจสอบกับของจริงในรีโปแล้ว
\newtcolorbox{goodbox}[1][]{guidebox, colback=goodwash, colframe=feg,
  colbacktitle=goodwash, coltitle=feg!65!black, #1}
% แดง — วิธีที่ผิดจริง ๆ อย่าทำ
\newtcolorbox{badbox}[1][]{guidebox, colback=badwash, colframe=bad,
  colbacktitle=badwash, coltitle=bad, #1}

% ---------------------------------------------------------------------
%  SEMANTIC MARKUP
% ---------------------------------------------------------------------
% ศัพท์อังกฤษที่แทรกในประโยคภาษาไทย ฟอนต์เนื้อความครอบคลุมละตินอยู่แล้วจึงไม่ต้อง
% สลับฟอนต์ — เก็บแมโครไว้เป็น "เครื่องหมายเชิงความหมาย" เผื่อวันหลังอยากเปลี่ยน
% การแสดงผลทั้งเล่มพร้อมกัน
\newcommand{\en}[1]{#1}
% ชื่อไฟล์ ฟังก์ชัน หรือค่าที่ปรากฏในโค้ดจริง
% ไม่บังคับขนาดตัวอักษร ปล่อยให้รับขนาดจากบริบท (สำคัญในตารางที่ใช้ \footnotesize
% มิฉะนั้นโค้ดจะกลายเป็นตัวใหญ่กว่าข้อความรอบข้างแล้วดันคอลัมน์ล้น)
% ฟอนต์ mono ถูกย่อไว้ 0.84 แล้วใน \setmonofont จึงพอดีกับเนื้อความอยู่แล้ว
\newcommand{\code}[1]{{\ttfamily #1}}

% ชื่อไฟล์ที่มีอักษรไทยอยู่ในชื่อ — ห้ามใช้ \code เพราะ DejaVu Sans Mono ไม่มี
% สัญลักษณ์ไทย ตัวอักษรจะหายไปจาก PDF โดยขึ้นแค่ warning "Missing character"
% ใช้ Garuda (sans ที่มีทั้งไทยและละติน) แทน
\newcommand{\fnth}[1]{{\sffamily\small #1}}

% ป้ายชื่อแนวคิด — สีเดียวกับที่ใช้ในไดอะแกรมและตาราง
\newcommand{\stg}[2]{\textcolor{#1}{\textbf{\en{#2}}}}

% ---------------------------------------------------------------------
%  ไวยากรณ์ของไดอะแกรม — นิยามครั้งเดียว ทุกรูปจะดูเป็นชุดเดียวกัน
%  หมายเหตุ: node ที่มีข้อความไทยต้องเผื่อ inner ysep มากกว่าไดอะแกรมภาษาละติน
%  เพราะสระบน/ล่างกินที่แนวตั้ง
% ---------------------------------------------------------------------
\tikzset{
  blk/.style={draw=rule, line width=0.6pt, rounded corners=2.5pt,
              fill=wash, minimum height=9mm, inner xsep=6pt, inner ysep=5pt,
              align=center, font=\sansall\footnotesize},
  stage/.style={blk, fill=ctr!8, draw=ctr!45},
  bestage/.style={blk, fill=beo!8, draw=beo!45},
  festage/.style={blk, fill=feg!8, draw=feg!45},
  opstage/.style={blk, fill=ops!8, draw=ops!45},
  ar/.style={-{Stealth[length=2.2mm,width=1.8mm]}, line width=0.7pt,
             draw=muted},
  lbl/.style={font=\sansall\scriptsize, text=muted, align=center},
  tag/.style={font=\sansall\scriptsize\bfseries, align=center},
}

% ---------------------------------------------------------------------
%  ตาราง
%  คอลัมน์แบบชิดซ้าย: ข้อความจัดชิดขอบสองข้าง (justified) ในคอลัมน์แคบจะเปิด
%  ช่องไฟกว้างมาก และแย่เป็นพิเศษกับภาษาไทย เพราะ ICU ตัดไทยเป็นก้อนที่ใหญ่กว่า
%  คำภาษาอังกฤษ → ใช้ L{} แทน p{} ในทุกคอลัมน์ที่มีข้อความยาว
% ---------------------------------------------------------------------
\newcolumntype{L}[1]{>{\raggedright\arraybackslash}p{#1}}
\captionsetup{skip=6pt}

% ---------------------------------------------------------------------
%  บล็อกโค้ด
%
%  ข้อจำกัดที่ต้องรู้: listings อ่านไฟล์แบบไบต์ ไม่เข้าใจ UTF-8 หลายไบต์
%  → "ห้ามใส่ภาษาไทยในบล็อกโค้ด" คอมเมนต์ในโค้ดทุกอันเป็นภาษาอังกฤษ
%    คำอธิบายภาษาไทยให้อยู่นอกบล็อกเสมอ
% ---------------------------------------------------------------------
\lstdefinelanguage{TypeScript}{
  keywords={await, async, break, case, catch, class, const, continue, default,
            delete, do, else, enum, export, extends, false, finally, for, from,
            function, if, implements, import, in, instanceof, interface, let,
            new, null, of, private, public, readonly, return, super, switch,
            this, throw, true, try, type, typeof, undefined, var, void, while,
            yield, satisfies, as},
  keywordstyle=\color{ctr}\bfseries,
  ndkeywords={string, number, boolean, Date, Promise, z, Record, Array},
  ndkeywordstyle=\color{ops},
  sensitive=true,
  comment=[l]{//},
  morecomment=[s]{/*}{*/},
  commentstyle=\color{muted}\itshape,
  stringstyle=\color{feg!70!black},
  morestring=[b]',
  morestring=[b]",
  morestring=[b]`,
}

\lstset{
  language=TypeScript,
  basicstyle=\ttfamily\scriptsize,
  backgroundcolor=\color{codebg},
  frame=leftline,
  framerule=2pt,
  rulecolor=\color{rule},
  framesep=7pt,
  xleftmargin=10pt,
  breaklines=true,
  breakatwhitespace=true,
  columns=fullflexible,
  keepspaces=true,
  showstringspaces=false,
  aboveskip=0.9em,
  belowskip=0.9em,
  literate={->}{{$\rightarrow$}}2 {=>}{{=>}}2,
}

% ชื่อ float ภาษาไทย — ถ้าไม่ตั้ง จะได้ "Figure 3" อยู่ใต้เนื้อความภาษาไทย
\renewcommand{\figurename}{รูปที่}
\renewcommand{\tablename}{ตารางที่}
\renewcommand{\refname}{เอกสารอ้างอิง}
\renewcommand{\contentsname}{สารบัญ}
% หมายเหตุ: ห้ามใส่ \color ใน \contentsname เพราะ \tableofcontents ส่งมันผ่าน
% \MakeUppercase ซึ่งจะเปลี่ยน ctr เป็น CTR แล้ว xcolor จะฟ้องว่าไม่รู้จักสีนั้น
% ถ้าจะจัดสไตล์หัวสารบัญ ให้ทำผ่าน \titleformat{\section} แทน
:
%    1) \hypersetup{pdftitle={...}}
%    2) \begin{document}
%
%  Build:  cd docs && latexmk trpc-guide.tex     (.latexmkrc บังคับ xelatex ให้แล้ว)
%          PDF ออกที่ docs/build/ แล้วค่อยคัดลอกไป docs/report/
% =============================================================================

\documentclass[11pt,a4paper]{article}

% ---------------------------------------------------------------------
%  FONTS
%
%  ต้องใช้ XeLaTeX เท่านั้น — pdflatex เรียงพิมพ์ภาษาไทยไม่ได้เลย (ไม่ใช่ "ได้แต่ไม่สวย")
%
%  ใช้ตระกูล TLWG ที่มีทั้งอักษรไทยและละตินอยู่ในไฟล์เดียว จงใจเลี่ยงการจับคู่
%  ฟอนต์ละตินกับฟอนต์ไทยล้วนผ่าน ucharclasses เพราะ ucharclasses สลับฟอนต์ตาม
%  "บล็อกยูนิโคด" ของตัวอักษร โดยไม่รู้จักขอบเขตของ TeX group → ข้อความละตินที่
%  ตามหลัง \textbf{...} จะยังค้างอยู่ในฟอนต์ไทยซึ่งไม่มีสัญลักษณ์ละติน แล้ว
%  "หายไปจาก PDF เงียบ ๆ" โดยไม่มี warning
% ---------------------------------------------------------------------
\usepackage{fontspec}

\setmainfont{Laksaman}[Ligatures=TeX]          % ไทย+ละติน — เนื้อความ
\setsansfont{Garuda}[Ligatures=TeX]            % ไทย+ละติน หนักกว่า — หัวข้อ/ป้ายในไดอะแกรม
\setmonofont{DejaVu Sans Mono}[Scale=0.84]     % 0.84 ทำให้ x-height ใกล้เคียง TLWG

% --- การตัดบรรทัดภาษาไทย ---
% ภาษาไทยเขียนติดกันไม่มีช่องว่าง TeX จึงหาจุดตัดบรรทัดเองไม่ได้ และจะปล่อยให้
% บรรทัดล้นออกนอกหน้ากระดาษ สองบรรทัดนี้ยกงานให้ ICU ตัดคำแบบพจนานุกรม
% ค่า stretch ต้องมากกว่าศูนย์อย่างมีนัยสำคัญ ไม่งั้น TeX แทบไม่มีอิสระในการจัด
% บรรทัดไทยและจะออกมาเป็น overfull box แทนที่จะยืด/หดช่องไฟ
\XeTeXlinebreaklocale "th"
\XeTeXlinebreakskip = 0pt plus 0.15em minus 0.02em
\emergencystretch=2.2em

% หัวข้อ ป้ายในไดอะแกรม หัวตาราง header/footer ใช้ sans — เนื้อความไม่ใช้
% นี่คือสิ่งที่แยก "โครงสร้างสำหรับกวาดสายตา" ออกจาก "เนื้อหาสำหรับอ่าน"
\newcommand{\sansall}{\sffamily}

% ---------------------------------------------------------------------
%  LAYOUT
% ---------------------------------------------------------------------
\usepackage[a4paper,margin=2.3cm,top=2.5cm,bottom=2.5cm,headsep=12pt]{geometry}
\usepackage{amsmath,amssymb}
\usepackage{graphicx}
\graphicspath{{figures/}}
\usepackage{booktabs}
\usepackage{array}
\usepackage{longtable}
\usepackage{multirow}
\usepackage{xcolor}
\usepackage{tikz}
\usetikzlibrary{arrows.meta,positioning,calc,shapes.geometric,shapes.misc,
                fit,backgrounds,decorations.pathmorphing}
\usepackage[most]{tcolorbox}
\usepackage[font=small,labelfont=bf,width=0.95\textwidth]{caption}
\usepackage{fancyhdr}
\usepackage{enumitem}
\usepackage{titlesec}
\usepackage{float}        % figure[H] — ตรึงรูปไว้ข้างเนื้อความที่อธิบายมัน
\usepackage{microtype}
\usepackage{listings}
\usepackage[colorlinks=true,allcolors=ctr,breaklinks=true,
            bookmarksnumbered=true]{hyperref}

% ภาษาไทยต้องการระยะบรรทัดมากกว่าละติน สระบนซ้อนวรรณยุกต์ และมีสระล่างด้วย
% ที่ระยะปกติจะชนกันข้ามบรรทัด — 1.28 คือค่าที่ทดสอบแล้ว
\linespread{1.28}
\setlength{\parskip}{0.35em}
\setlength{\parindent}{0pt}

% ---------------------------------------------------------------------
%  สี — ระบบสีเชิงความหมาย (semantic colour)
%
%  ฐานสีคือชุด Okabe–Ito ซึ่งปลอดภัยสำหรับผู้ที่มองสีบางคู่ไม่แยก (colour-blind safe)
%
%  สีสี่สีหลักถูก "ตั้งชื่อกลาง ๆ" ไว้ที่นี่ แล้วให้แต่ละเอกสารประกาศความหมายของ
%  ตัวเองในหน้า "เอกสารนี้อ่านอย่างไร":
%
%    trpc-guide.tex   → 4 ชั้นของระบบ:   สัญญา / backend / frontend / เครือข่าย
%    trpc-meeting.tex → 4 ระดับเร่งด่วน: P0 / P1 / P2 / P3
%
%  สองเอกสารใช้คนละความหมายได้ เพราะแต่ละฉบับ "ประกาศระบบสีของตัวเองไว้หน้าแรก"
%  และคงเส้นคงวาภายในเล่ม — สิ่งที่ห้ามคือเปลี่ยนความหมายกลางเล่ม
% ---------------------------------------------------------------------
\definecolor{ctr}{HTML}{0072B2}      % ฟ้า   — สัญญา/schema/type   · หรือ P2
\definecolor{beo}{HTML}{D55E00}      % ส้ม   — ฝั่ง backend         · หรือ P1
\definecolor{feg}{HTML}{009E73}      % เขียว — ฝั่ง frontend        · หรือ P3
\definecolor{ops}{HTML}{7B5EA7}      % ม่วง  — เครือข่าย/รันไทม์/ops

\definecolor{ink}{HTML}{1A2733}      % เข้ม — หัวข้อรอง
\definecolor{muted}{HTML}{5B6B7A}    % ข้อความรอง เส้นคั่น ป้ายกำกับ
\definecolor{rule}{HTML}{C9D2DB}     % เส้นบาง ขอบไดอะแกรม
\definecolor{wash}{HTML}{F1F6FA}     % พื้นกล่องฟ้า/กลาง
\definecolor{warnwash}{HTML}{FDF3E7} % พื้นกล่องส้ม
\definecolor{goodwash}{HTML}{E9F4F0} % พื้นกล่องเขียว
\definecolor{badwash}{HTML}{FBEDED}  % พื้นกล่องแดง
\definecolor{bad}{HTML}{B3261E}      % แดง — ผิด/อันตราย · หรือ P0
\definecolor{codebg}{HTML}{F7F9FB}   % พื้นบล็อกโค้ด

% ---------------------------------------------------------------------
%  HEADINGS
%  \part ใช้เลขอารบิก เพื่อให้อ้างอิงข้ามเล่มอ่านว่า "ภาค 3" ไม่ใช่ "ภาค III"
% ---------------------------------------------------------------------
\renewcommand{\thepart}{\arabic{part}}
\titleformat{\part}[display]
  {\sansall\huge\bfseries\color{ctr}}
  {\normalsize\color{muted}ภาค \thepart}{0.4em}{}
\titlespacing*{\part}{0pt}{0pt}{1.5em}

\titleformat{\section}
  {\sansall\Large\bfseries\color{ctr!88!black}}{\thesection}{0.7em}{}
\titleformat{\subsection}
  {\sansall\large\bfseries\color{ink}}{\thesubsection}{0.6em}{}
\titleformat{\subsubsection}
  {\sansall\normalsize\bfseries\color{muted}}{\thesubsubsection}{0.55em}{}
\titlespacing*{\section}{0pt}{1.9em}{0.7em}
\titlespacing*{\subsection}{0pt}{1.3em}{0.45em}
\titlespacing*{\subsubsection}{0pt}{1.0em}{0.35em}

% หัวกระดาษ: ชื่อเอกสารคงที่ทางซ้าย ชื่อหัวข้อปัจจุบันทางขวา ผู้อ่านจะรู้ตลอดว่าอยู่ตรงไหน
% \docshort ถูกกำหนดในไฟล์เอกสารแต่ละฉบับ
\providecommand{\docshort}{}
\pagestyle{fancy}
\fancyhf{}
\renewcommand{\headrulewidth}{0.4pt}
\fancyhead[L]{\footnotesize\sansall\color{muted}\docshort}
\fancyhead[R]{\footnotesize\sansall\color{muted}\leftmark}
\fancyfoot[C]{\footnotesize\color{muted}\thepage}
\renewcommand{\sectionmark}[1]{\markboth{#1}{}}

% รายการแบบกระชับ — ค่าเริ่มต้นของ LaTeX ดูหลวมเมื่ออยู่กับ \parskip
\setlist[itemize]{itemsep=1pt,topsep=3pt,parsep=0pt,leftmargin=1.4em}
\setlist[enumerate]{itemsep=1pt,topsep=3pt,parsep=0pt,leftmargin=1.6em}
\setlist[description]{itemsep=2pt,topsep=3pt,parsep=0pt}

% ---------------------------------------------------------------------
%  กล่องข้อความ 4 แบบ ความหมายตายตัว ห้ามคิดแบบที่ห้ากลางเล่ม
%
%  colbacktitle ต้องตั้งให้เท่ากับ colback มิฉะนั้น tcolorbox จะใช้ colframe
%  (สีเข้ม) เป็นพื้นหลังของหัวกล่อง แล้วตัวอักษรหัวกล่องสีเข้มจะอ่านไม่ออก
% ---------------------------------------------------------------------
\tcbset{guidebox/.style={
  enhanced, breakable, boxrule=0pt, leftrule=2.5pt, arc=1pt,
  left=8pt, right=8pt, top=6pt, bottom=6pt,
  fonttitle=\sansall\bfseries\small, titlerule=0pt,
  toptitle=1pt, bottomtitle=1pt}}

% ฟ้า — ใจความที่ต้องจำ / ประโยคที่ร้อยหลายบทเข้าด้วยกัน
\newtcolorbox{keybox}[1][]{guidebox, colback=wash, colframe=ctr,
  colbacktitle=wash, coltitle=ctr!85!black, #1}
% ส้ม — กับดัก หรือข้อสมมติที่ถ้าลืมแล้วข้อสรุปจะผิด
\newtcolorbox{warnbox}[1][]{guidebox, colback=warnwash, colframe=beo,
  colbacktitle=warnwash, coltitle=beo!80!black, #1}
% เขียว — สิ่งที่ตรวจสอบกับของจริงในรีโปแล้ว
\newtcolorbox{goodbox}[1][]{guidebox, colback=goodwash, colframe=feg,
  colbacktitle=goodwash, coltitle=feg!65!black, #1}
% แดง — วิธีที่ผิดจริง ๆ อย่าทำ
\newtcolorbox{badbox}[1][]{guidebox, colback=badwash, colframe=bad,
  colbacktitle=badwash, coltitle=bad, #1}

% ---------------------------------------------------------------------
%  SEMANTIC MARKUP
% ---------------------------------------------------------------------
% ศัพท์อังกฤษที่แทรกในประโยคภาษาไทย ฟอนต์เนื้อความครอบคลุมละตินอยู่แล้วจึงไม่ต้อง
% สลับฟอนต์ — เก็บแมโครไว้เป็น "เครื่องหมายเชิงความหมาย" เผื่อวันหลังอยากเปลี่ยน
% การแสดงผลทั้งเล่มพร้อมกัน
\newcommand{\en}[1]{#1}
% ชื่อไฟล์ ฟังก์ชัน หรือค่าที่ปรากฏในโค้ดจริง
% ไม่บังคับขนาดตัวอักษร ปล่อยให้รับขนาดจากบริบท (สำคัญในตารางที่ใช้ \footnotesize
% มิฉะนั้นโค้ดจะกลายเป็นตัวใหญ่กว่าข้อความรอบข้างแล้วดันคอลัมน์ล้น)
% ฟอนต์ mono ถูกย่อไว้ 0.84 แล้วใน \setmonofont จึงพอดีกับเนื้อความอยู่แล้ว
\newcommand{\code}[1]{{\ttfamily #1}}

% ชื่อไฟล์ที่มีอักษรไทยอยู่ในชื่อ — ห้ามใช้ \code เพราะ DejaVu Sans Mono ไม่มี
% สัญลักษณ์ไทย ตัวอักษรจะหายไปจาก PDF โดยขึ้นแค่ warning "Missing character"
% ใช้ Garuda (sans ที่มีทั้งไทยและละติน) แทน
\newcommand{\fnth}[1]{{\sffamily\small #1}}

% ป้ายชื่อแนวคิด — สีเดียวกับที่ใช้ในไดอะแกรมและตาราง
\newcommand{\stg}[2]{\textcolor{#1}{\textbf{\en{#2}}}}

% ---------------------------------------------------------------------
%  ไวยากรณ์ของไดอะแกรม — นิยามครั้งเดียว ทุกรูปจะดูเป็นชุดเดียวกัน
%  หมายเหตุ: node ที่มีข้อความไทยต้องเผื่อ inner ysep มากกว่าไดอะแกรมภาษาละติน
%  เพราะสระบน/ล่างกินที่แนวตั้ง
% ---------------------------------------------------------------------
\tikzset{
  blk/.style={draw=rule, line width=0.6pt, rounded corners=2.5pt,
              fill=wash, minimum height=9mm, inner xsep=6pt, inner ysep=5pt,
              align=center, font=\sansall\footnotesize},
  stage/.style={blk, fill=ctr!8, draw=ctr!45},
  bestage/.style={blk, fill=beo!8, draw=beo!45},
  festage/.style={blk, fill=feg!8, draw=feg!45},
  opstage/.style={blk, fill=ops!8, draw=ops!45},
  ar/.style={-{Stealth[length=2.2mm,width=1.8mm]}, line width=0.7pt,
             draw=muted},
  lbl/.style={font=\sansall\scriptsize, text=muted, align=center},
  tag/.style={font=\sansall\scriptsize\bfseries, align=center},
}

% ---------------------------------------------------------------------
%  ตาราง
%  คอลัมน์แบบชิดซ้าย: ข้อความจัดชิดขอบสองข้าง (justified) ในคอลัมน์แคบจะเปิด
%  ช่องไฟกว้างมาก และแย่เป็นพิเศษกับภาษาไทย เพราะ ICU ตัดไทยเป็นก้อนที่ใหญ่กว่า
%  คำภาษาอังกฤษ → ใช้ L{} แทน p{} ในทุกคอลัมน์ที่มีข้อความยาว
% ---------------------------------------------------------------------
\newcolumntype{L}[1]{>{\raggedright\arraybackslash}p{#1}}
\captionsetup{skip=6pt}

% ---------------------------------------------------------------------
%  บล็อกโค้ด
%
%  ข้อจำกัดที่ต้องรู้: listings อ่านไฟล์แบบไบต์ ไม่เข้าใจ UTF-8 หลายไบต์
%  → "ห้ามใส่ภาษาไทยในบล็อกโค้ด" คอมเมนต์ในโค้ดทุกอันเป็นภาษาอังกฤษ
%    คำอธิบายภาษาไทยให้อยู่นอกบล็อกเสมอ
% ---------------------------------------------------------------------
\lstdefinelanguage{TypeScript}{
  keywords={await, async, break, case, catch, class, const, continue, default,
            delete, do, else, enum, export, extends, false, finally, for, from,
            function, if, implements, import, in, instanceof, interface, let,
            new, null, of, private, public, readonly, return, super, switch,
            this, throw, true, try, type, typeof, undefined, var, void, while,
            yield, satisfies, as},
  keywordstyle=\color{ctr}\bfseries,
  ndkeywords={string, number, boolean, Date, Promise, z, Record, Array},
  ndkeywordstyle=\color{ops},
  sensitive=true,
  comment=[l]{//},
  morecomment=[s]{/*}{*/},
  commentstyle=\color{muted}\itshape,
  stringstyle=\color{feg!70!black},
  morestring=[b]',
  morestring=[b]",
  morestring=[b]`,
}

\lstset{
  language=TypeScript,
  basicstyle=\ttfamily\scriptsize,
  backgroundcolor=\color{codebg},
  frame=leftline,
  framerule=2pt,
  rulecolor=\color{rule},
  framesep=7pt,
  xleftmargin=10pt,
  breaklines=true,
  breakatwhitespace=true,
  columns=fullflexible,
  keepspaces=true,
  showstringspaces=false,
  aboveskip=0.9em,
  belowskip=0.9em,
  literate={->}{{$\rightarrow$}}2 {=>}{{=>}}2,
}

% ชื่อ float ภาษาไทย — ถ้าไม่ตั้ง จะได้ "Figure 3" อยู่ใต้เนื้อความภาษาไทย
\renewcommand{\figurename}{รูปที่}
\renewcommand{\tablename}{ตารางที่}
\renewcommand{\refname}{เอกสารอ้างอิง}
\renewcommand{\contentsname}{สารบัญ}
% หมายเหตุ: ห้ามใส่ \color ใน \contentsname เพราะ \tableofcontents ส่งมันผ่าน
% \MakeUppercase ซึ่งจะเปลี่ยน ctr เป็น CTR แล้ว xcolor จะฟ้องว่าไม่รู้จักสีนั้น
% ถ้าจะจัดสไตล์หัวสารบัญ ให้ทำผ่าน \titleformat{\section} แทน
